\documentclass[a4paper,12pt]{article}
\usepackage{color}
\usepackage{graphicx, gensymb}
\usepackage{amsfonts, cancel, amssymb}
\usepackage{amsmath, amsthm, array, esvect, esint}
\usepackage{typearea}
\usepackage[a4paper,left=2cm,right=2cm,top=2cm,bottom=3cm,bindingoffset=5mm]{geometry}
\newcommand\numberthis{\addtocounter{equation}{1}\tag{\theequation}}
\newcommand{\bra}{}
\newcommand{\kom}{}
\newcommand{\brak}{}
\newcommand{\braket}{}
\newcommand{\intx}{}
\newcommand{\intr}{}
\newcommand{\kla}{}
\newcommand{\klam}{}
\newcommand{\klammer}{}
\newcommand{\oper}{}
\newcommand{\tecxt}{}
\newcommand{\Bra}{}
\newcommand{\Ket}{}

\begin{document}

\title{9 Ferromagnetismus}
\author{Joachim Schwardt}
\date{\today}
\maketitle

\section*{Aufgabe 9.1: Ferromagnetismus}
Betrachte ein $2\times L$-Gitter von Spins unter periodischen Randbedingungen. Dann lautet der Hamiltonian
\begin{align*}
H&=-J\sum_{j=1}^2\sum_{k=1}^L\sigma_{j,k}\sigma_{j,k+1}-2J\sum_{k=1}^L\sigma_{j,k}\sigma_{j+1,k}
\end{align*}
Betrachte eine Zelle mit vier Gitterplätzen. Es gibt $2^4=16$ mögliche Zustände, die sich in der Transfermatrix zusammenfassen lassen:
\begin{align*}
H_{k,k+1}&=-J\klam\left[ {\sigma_{1,k}\sigma_{1,k+1}+\sigma_{2,k}\sigma_{2,k+1}+2\sigma_{1,k}\sigma_{2,k}} \right]\\
T_{k,k+1}&=\begin{pmatrix}
e^{-\beta JH(\sigma_{1k},\sigma_{2k},\sigma_{1k+1},\sigma_{2k+1}}
\end{pmatrix}=\\
&=\begin{pmatrix}
(++++)&(+++-)&\dots\\
(+-++)&\dots\\
(-+++)&\dots\\
(--++)&\dots
\end{pmatrix} =\begin{pmatrix}
e^{4\beta J}&e^{2\beta J}&e^{2\beta J}&1\\
e^{-2\beta J}&1&e^{-4\beta J}&e^{-2\beta J}\\
e^{-2\beta J}&e^{-4\beta J}&1&e^{-2\beta J}\\
1&e^{2\beta J}&e^{2\beta J}&e^{4\beta J}
\end{pmatrix}
\end{align*}
Die Berechnung der Eigenwerte erfolgt mit $z=e^{2\beta J}$:
\begin{align*}
\det (T-\lambda)&=\det\begin{pmatrix}
z^2-\lambda&z&z&1\\
\frac{1}{z}&1-\lambda&\frac{1}{z^2}&\frac{1}{z}\\
\frac{1}{z}&\frac{1}{z^2}&1-\lambda&\frac{1}{z}\\
1&z&z&z^2-\lambda
\end{pmatrix}=\\
&=\det\begin{pmatrix}
z^2-1-\lambda&z&z&1\\
0&1-\lambda&\frac{1}{z^2}&\frac{1}{z}\\
0&\frac{1}{z^2}&1-\lambda&\frac{1}{z}\\
1-z^2+\lambda&z&z&z^2-\lambda
\end{pmatrix}=\\
&=\det\begin{pmatrix}
z^2-1-\lambda&z&z&1\\
0&1-\lambda&\frac{1}{z^2}&\frac{1}{z}\\
0&\frac{1}{z^2}&1-\lambda&\frac{1}{z}\\
0&2z&2z&z^2+1-\lambda
\end{pmatrix}=\\
&=\det\begin{pmatrix}
z^2-1-\lambda&0&z&1\\
0&1-\frac{1}{z^2}-\lambda&\frac{1}{z^2}&\frac{1}{z}\\
0&0&1+\frac{1}{z^2}-\lambda&\frac{2}{z}\\
0&0&2z&z^2+1-\lambda
\end{pmatrix}=\\
&=\kla\left( {z^2-1-\lambda} \right)\kla\left( {1-\frac{1}{z^2}-\lambda} \right)\klam\left[ {\kla\left( {1+\frac{1}{z^2}-\lambda} \right)(z^2+1-\lambda)-4} \right]
\end{align*}
Setze nun $z^2=e^{4\beta J}\rightarrow z$. Die ersten beiden Eigenwerte lautet dann $\lambda_1=z-1$ und $\lambda_2=1-\frac{1}{z}$. Gesucht sind die restlichen Nullstellen des Polynoms:
\begin{align*}
0&\overset{!}{=}\lambda^2-\lambda\kla\left( {z+2+\frac{1}{z}} \right)-4+\kla\left( {1+\frac{1}{z}} \right)\kla\left( {1+z} \right)\\
\Rightarrow\ \ \ \lambda_{3,4}&=\frac{z+2+z^{-1}}{2}\pm\frac{\sqrt{z^2+4+z^{-2}+4z+4z^{-1}+2-4\kla\left( {z+2+ z^{-1}-4} \right)}}{2}=\\
&=\frac{z+2+z^{-1}}{2}\pm\frac{\sqrt{z^2+14+z^{-2}}}{2}
\end{align*}
Es ist $z>1$, woraus $\lambda_1=z\lambda_2>\lambda_2$ folgt. Es ist offensichtlich $\lambda_3>\lambda_4$. Insgesamt gilt $\lambda_3\ge\lambda_j$:
\begin{align*}
\lambda_3&\ge\frac{z+2}{2}+\frac{\sqrt{z^2}}{2}=z+1>\lambda_{1,2}
\end{align*}
Die Zustandssumme lautet:
\begin{align*}
Z&=\sum_\sigma\exp\kla\left( {-\beta H} \right)=\sum_\sigma\prod_k\exp\kla\left( {\beta J\klam\left[ {\sigma_{1,k}\sigma_{1,k+1}+\sigma_{2,k}\sigma_{2,k+1}+2\sigma_{1,k}\sigma_{2,k}} \right]} \right)=\\
&=\sum_{\sigma_{11},\sigma_{21}}\dots\sum_{\sigma_{1L},\sigma_{2L}}\braket\langle {\sigma_{11},\sigma_{21}} | {T} | {\dots}\rangle\dots\braket\langle {\sigma_{1L},\sigma_{2L}} | {T} | {\sigma_{11},\sigma_{21}}\rangle=\\
&=\sum_{\sigma_{11},\sigma_{21}}\braket\langle {\sigma_{11},\sigma_{21}} | {T^L} | {\sigma_{11},\sigma_{21}}\rangle=\tecxt\text{tr\,}T^L=\lambda_1^L+\lambda_2^L+\lambda_3^L+\lambda_4^L
\end{align*}
Im Grenzwert großer Systeme $L\rightarrow\infty$ ist nur der größte Eigenwert relevant:
\begin{align*}
Z&\sim\lambda_3^L=\frac{1}{2^L}\klam\left[ {z+2+\frac{1}{z}+\sqrt{z^2+14+\frac{1}{z^2}}} \right]^L
\end{align*}
Betrachte nun die mittlere freie Energie mit $F=-\frac{1}{\beta}\log Z$:
\begin{align*}
f&=\lim\frac{F}{L}=\frac{\log 2}{\beta}-\frac{1}{\beta}\log\klam\left[ {z+2+\frac{1}{z}+\sqrt{z^2+14+\frac{1}{z^2}}} \right]
\end{align*}
\section*{Aufgabe 9.2: Landau-Theorie des 2D Ising Modells}


\end{document}